\subsection{Applications: III. Algebras of continuous functions}

Lemma 4. --- Let X be a compact space, H a closed linear subspace of the Banach space \(\mathcal{C}(\mathrm{X};\mathbf{C})\) (resp. \(\mathcal{C}(\mathrm{X};\mathbf{R})\) ). Let a be a point of X admitting a countable fundamental system of neighborhoods; assume that, for any numbers c and d such that 0 < c < d < 1 and any open neighborhood U of a, there exists an \(f \in H\) such that

\[
| f | \leqslant 1, \quad | f (a) | \geqslant d, \quad | f (x) | \leqslant c \text {   for   all   } x \in X - U.\tag{12}
\]

Then there exists a function \(u \in \mathcal{H}\) such that \(|u(x)| < |u(a)|\) for all \(x \neq a\) .

Let \(\left(\mathrm{V}_{n}\right)\) ( \(n \geqslant 1\) ) be a fundamental system of neighborhoods of \(a\) , and let \(\lambda, \mu, \varepsilon\) be numbers such that

\[
0 <   \lambda <   1, \quad 1 <   \mu <   \mu + \varepsilon \leqslant 1 + \lambda .
\]

Thus \(0 < \lambda / \mu < 1 / \mu < 1\) . We are going to define, by induction on \(n\) ( \(n \geqslant 1\) ), a decreasing sequence \((\mathrm{U}_n)\) of open neighborhoods of \(a\) such that \(\mathrm{U}_n \subset \mathrm{V}_n\) for all \(n\) , and a sequence \((h_n)\) of functions in \(\mathcal{H}\) satisfying the relations

\[
\left| h _ {n} (x) \right| \leqslant \mu \quad \text { for   all } x \in X\tag{\ensuremath{13_n}}
\]

\[
h _ {n} (a) = 1\tag{\ensuremath{14_n}}
\]

\[
\left| h _ {n} (x) \right| \leqslant \lambda \quad \text { for   all } x \in X - U _ {n}\tag{\ensuremath{15_n}}
\]

\[
\left| \sum_ {j = 1} ^ {n} \lambda^ {j} h _ {j} (y) \right| <   \sum_ {j = 1} ^ {n + 1} \lambda^ {j} \quad \text { for   all } y \in X.\tag{\ensuremath{16_n}}
\]

Assume \(h_m\) and \(U_m\) defined for \(1 \leqslant m < n\) , satisfying the four preceding conditions (with \(n\) replaced by \(m\) ); on the other hand, set \(U_0 = X\) . The function \(\sum_{j=1}^{n-1} \lambda^j h_j\) (equal to 0 if \(n = 1\) ) is continuous and takes the value \(\sum_{j=1}^{n-1} \lambda^j\) at the point \(a\) ; therefore there exists an open neighborhood \(U_n\) of \(a\) , contained in \(U_{n-1} \cap V_n\) , such that

\[
\left| \sum_ {j = 1} ^ {n - 1} \lambda^ {j} h _ {j} (y) \right| <   \sum_ {j = 1} ^ {n - 1} \lambda^ {j} + \varepsilon \lambda^ {n} \quad \text { for   all } y \in \mathrm{U} _ {n}.\tag{17}
\]

By hypothesis there exists a function \(f \in \mathcal{H}\) such that

\[
\begin{array}{l} | f (x) | \leqslant 1 \quad \text { for   all } x \in X, \quad | f (a) | \geqslant 1 / \mu , \\ | f (x) | \leqslant \lambda / \mu \quad \text { for } x \in X - U _ {n}. \end{array}
\]

Set \(h_n = f / f(a)\) ; the relations \((13_n), (14_n)\) and \((15_n)\) then hold; set

\[
g = \sum_ {j = 1} ^ {n} \lambda^ {j} h _ {j} = \sum_ {j = 1} ^ {n - 1} \lambda^ {j} h _ {j} + \lambda^ {n} h _ {n}.
\]

By (17) and \((13_{n})\) , for \(y \in \mathrm{U}_n\) we have

\[
| g (y) | <   \sum_ {j = 1} ^ {n - 1} \lambda^ {j} + \varepsilon \lambda^ {n} + \mu \lambda^ {n} \leqslant \sum_ {j = 1} ^ {n + 1} \lambda^ {j},
\]

since \(\varepsilon + \mu \leqslant 1 + \lambda\) ; for \(x \in \mathrm{X} - \mathrm{U}_n\) , we have \(|h_p(x)| \leqslant \lambda\) for \(1 \leqslant p \leqslant n\) , whence also

\[
| g (x) | \leqslant \sum_ {j = 2} ^ {n + 1} \lambda^ {j} <   \sum_ {j = 1} ^ {n + 1} \lambda^ {j},
\]

which completes the proof of \((16_{n})\) .

This being so, the series \(\sum_{n=1}^{\infty} \lambda^n h_n\) is normally convergent in X since \(\lambda < 1\) and \(|h_n(x)| \leqslant \mu\) for all \(n\) and all \(x \in \mathbf{X}\) ; let \(u\) be its sum, which belongs to \(\mathcal{H}\) since \(\mathcal{H}\) is closed. By the relation \((14_n)\) , we have \(u(a) = \sum_{n=1}^{\infty} \lambda^n\) ; on the other hand if \(x \neq a\) , there exists an integer \(n\) such that \(x \notin \mathrm{U}_{n+1}\) ; therefore \(|h_{n+k}(x)| \leqslant \lambda\) for all \(k \geqslant 1\) by the relation \((15_n)\) ; it follows, using \((16_n)\) , that

\[
\begin{array}{l} | u (x) | \leqslant \left| \sum_ {j = 1} ^ {n} \lambda^ {j} h _ {j} (x) \right| + \left| \sum_ {j = n + 1} ^ {\infty} \lambda^ {j} h _ {j} (x) \right| <   \sum_ {j = 1} ^ {n + 1} \lambda^ {j} + \lambda \sum_ {j = n + 1} ^ {\infty} \lambda^ {j} \\ = \sum_ {j = 1} ^ {\infty} \lambda^ {j} = | u (a) |. \end{array}
\]

THEOREM 2 (E. Bishop). --- Let X be a compact space, A a closed subalgebra of the complex Banach algebra \(\mathcal{C}(\mathrm{X};\mathbf{C})\) . Assume that A contains the constants and separates the points of X. Let \(a\) be a point of X; the following conditions are equivalent:

\begin{enumerate}
\def\labelenumi{\alph{enumi})}
\item
  There exists a function \(f \in \mathcal{A}\) such that \(|f(x)| < |f(a)|\) for all \(x \neq a\) .
\item
  The point \(a\) is \(\mathcal{A}\) -extremal and admits a countable fundamental system of neighborhoods.
\end{enumerate}

\(a) \Rightarrow b)\) : Let \(f \in \mathcal{A}\) be such that \(|f(a)| > |f(x)|\) for \(x \neq a\) ; by Prop. 9 of No.~4, \(a\) is an \(\mathcal{A}\) -extremal point. On the other hand, if \(\mathrm{U}_n\) is the set of \(x \in \mathrm{X}\) such that \(|f(x)| > |f(a)| - 1/n\) , then \(\mathrm{U}_n\) is an open neighborhood of \(a\) , and the intersection of the \(\mathrm{U}_n\) reduces to \(a\) ; since \(\mathrm{X}\) is compact, the \(\mathrm{U}_n\) form a fundamental system of neighborhoods of \(a\) (GT, I, §9, No.~1, Th. 1).

\(b) \Rightarrow a)\) : It suffices to verify that \(b)\) implies the hypotheses of Lemma 4. With the notations of that lemma, set \(\varepsilon = \log d / \log c\) ; thus \(0 < \varepsilon < 1\) . Since \(a\) is an \(\mathcal{A}_r\) -extremal point, there exists a function \(g \in \dot{\mathcal{A}}\) such that

\[
\mathcal {R} (g) \geqslant 0, \quad \mathcal {R} (g (a)) \leqslant \varepsilon , \quad \mathcal {R} (g (x)) \geqslant 1 \text {for} x \in X - U
\]

(No.~3, Prop. 6, b)). Set \(f = c^g\) ; since \(f\) is the sum of the normally convergent series \(\sum_{n=0}^{\infty} (\log c)^n g^n / n!\) , we have \(f \in \mathcal{A}\) and

\[
| f | \leqslant 1, \quad | f (a) | \geqslant c ^ {\varepsilon} = d, \quad | f (x) | \leqslant c \text { for } x \in X - U.\tag{Q.E.D.}
\]

COROLLARY. --- Suppose in addition that X is metrizable. Then the following properties are equivalent:

\begin{enumerate}
\def\labelenumi{\alph{enumi})}
\item
  \(a\) is an \(\mathcal{A}\) -extremal point of X.
\item
  There exists \(u \in \mathcal{A}\) such that \(|u(x)| < |u(a)|\) for all \(x \neq a\) .
\item
  Let \(\mathfrak{M}\) be the set of subsets \(M\) of \(X\) such that for every function \(f \in \mathcal{A}\) , \(|f|\) attains its supremum in \(X\) at at least one point of \(M\) . Then a belongs to all of the sets \(M \in \mathfrak{M}\) .
\item
  Let \(\mathfrak{N}\) be the set of subsets N of X such that, for every function \(f \in \mathcal{A}\) , \(\mathcal{R}(f)\) attains its supremum in X at at least one point of N. Then a belongs to all of the sets \(N \in \mathfrak{N}\) .
\end{enumerate}

In other words,

\[
\operatorname{Ch} _ {\mathcal {A}} (\mathrm{X}) = \bigcap_ {\mathrm{M} \in \mathfrak {M}} \mathrm{M} = \bigcap_ {\mathrm{N} \in \mathfrak {N}} \mathrm{N}.\tag{18}
\]

Since, in a metrizable space, every point admits a countable fundamental system of neighborhoods, the equivalence of \(a\) and \(b\) follows from Th. 2. Let us show that \(b\) implies \(c\) : indeed, \(a\) is the unique point where \(|u|\) attains its supremum; on the other hand, \(c\) implies \(a\) because, for every \(f \in \mathcal{A}\) , \(\mathrm{Ch}_{\mathcal{A}}(\mathrm{X})\) intersects the set of points where \(|f|\) attains its supremum (No.~4, Prop. 9). The same reasoning, using Prop. 5 of No.~3, shows that \(d\) implies \(a\) ). Finally, to see that \(b\) implies \(d\) , we can restrict ourselves to the case that X does not reduce to the single point \(a\) , therefore \(u(a) \neq 0\) ; the function \(v = u / u(a)\) then belongs to \(\mathcal{A}\) , and we have \(v(a) = 1\) and \(|v(x)| < 1\) for \(x \neq a\) , whence \(\mathcal{R}(v(a)) = 1\) and \(\mathcal{R}(v(x)) < 1\) for \(x \neq a\) . Since the function \(\mathcal{R}(v)\) attains its supremum only at the point \(a\) , we have indeed \(a \in N\) for every \(N \in \mathfrak{N}\) .

*Examples. --- Let \(\mathrm{X}_1\) be the set of points \((z_1, z_2) \in \mathbf{C}^2\) such that \(|z_1|^2 + |z_2|^2 \leqslant 1\) (the unit ball in \(\mathbf{R}^4\) ) and let \(\mathcal{A}_1'\) be the set of restrictions to \(\mathrm{X}_1\) of the holomorphic functions, with values in \(\mathbf{C}\) , defined in a neighborhood of \(\mathrm{X}_1\) in \(\mathbf{C}^2\) (the neighborhood depending on the function considered); let \(\mathcal{A}_1\) be the closure of \(\mathcal{A}_1'\) in \(\mathcal{C}(\mathrm{X}_1; \mathbf{C})\) , which is obviously a closed complex subalgebra of \(\mathcal{C}(\mathrm{X}_1; \mathbf{C})\) and separates the points of \(\mathrm{X}_1\) . Application of the `maximum principle' for holomorphic functions shows that \(\mathrm{Ch}_{\mathcal{A}_1}(\mathrm{X}_1)\) is the sphere \(\mathbf{S}_3\) .

In the preceding definition, let us replace \(X_1\) by the `polydisk' \(X_2\) defined by the relations \(|z_1| \leqslant 1\) and \(|z_2| \leqslant 1\) , which yields subalgebras \(\mathscr{A}_2'\) and \(\mathscr{A}_2\) (the closure of \(\mathcal{A}_2^{\prime}\) ) of \(\mathcal{C}(\mathrm{X}_2; \mathbf{C})\) . Here, the maximum principle shows that \(\mathrm{Ch}_{\mathcal{A}_2}(\mathrm{X}_2)\) is the `torus' defined by the relations \(|z_1| = 1\) and \(|z_2| = 1\) .

From these results, one deduces that there does not exist an analytic isomorphism of an open neighborhood of \(X_{1}\) onto an open neighborhood of \(X_{2}\) that transforms \(X_{1}\) into \(X_{2}\) ; for, if v were the restriction to \(X_{1}\) of such a mapping, one would have \(A_{2}=vA_{1}v^{-1}\) and so v would transform \(S_{3}\) into a space homeomorphic to \(T^{2}\) , which is absurd since \(S_{3}\) is simply connected but \(T^{2}\) is not. One will observe, however, that the spaces \(X_{1}\) and \(X_{2}\) are homeomorphic, both being bounded convex sets in \(R^{4}\) with nonempty interior.

